\documentclass[10pt,a4paper]{amsart}
\usepackage[margin=19mm]{geometry}
\usepackage{amsmath,amssymb,mathtools}
\usepackage[scr=rsfs,cal=euler]{mathalfa}
\usepackage{xfrac}
\usepackage[unicode,hidelinks,bookmarks=false]{hyperref}
\newcommand{\ie}{i.e.\ }
\newcommand{\cf}{cf.\ }
\newcommand{\resp}{respectively\ }
\newcommand{\Subsection}{\subsection}
\providecommand{\nobreakdash}{\nobreak\hbox{-}}
\setlength{\emergencystretch}{2em}
\multlinegap0pt
\usepackage{bm}
\usepackage{bbm}
\usepackage{dsfont}
\usepackage{xspace}
\usepackage{cancel}
\usepackage{mathtools}
\usepackage{amsthm}
\makeatletter
\@ifundefined{theorem}{\newtheorem{theorem}{Theorem}}{}
\@ifundefined{lemma}{\newtheorem{lemma}[theorem]{Lemma}}{}
\makeatother
\newcommand{\bbx}{\bm x}
\newcommand{\rd}{\,\mathrm{d}}
\title[Fully discrete energy-dissipative and conservative discrete ]{Fully discrete energy-dissipative and conservative discrete gradient particle methods for a class of continuity equations}
\author{Jingwei Hu, Samuel Q. Van Fleet, Andy T. S. Wan}
\date{Source: arXiv:2407.00533v1 (CC BY 4.0)}
\begin{document}
\maketitle

\noindent\emph{Source and licence.} Fully discrete energy-dissipative and conservative discrete gradient particle methods for a class of continuity equations. Version arXiv:2407.00533v1 (CC BY 4.0), distributed under the Creative Commons Attribution 4.0 International licence, as recorded in the arXiv licence metadata for this exact version and shown on the abstract page https://arxiv.org/abs/2407.00533v1. Full rights notice: licenses/arxiv-2407.00533-CC-BY-4.0.txt.\par\smallskip

\noindent\textit{Editorial scope.} Counted unit: the Lemma whose source label is \texttt{lem:compCond} in the article (source0.tex lines 324--330, source label \texttt{lem:compCond}), delivered together with the article's own complete original proof of it (source0.tex lines 331--365, one \texttt{proof} environment of 35 lines). No mathematics is written, repaired or re-derived: statement and proof are the article's own text line for line. Required context retained uncounted: fparticle (lines 281--283); eq:regularE (lines 320--322). Every definition, display and label of those lines is retained unchanged. Omitted from this package and not redistributed: 1--280, 284--319, 323--323, 366--953. Nothing inside a retained excerpt is omitted or altered: every retained body is delivered exactly as the article's lines are written, so no omission receipt of any kind accompanies this package. Citations: the retained excerpts carry no citation command at all, so no bibliography entry of the article is transcribed and no reference is redistributed. Reference adaptation: none was needed and none was made; every reference inside the delivered unit resolves inside it, so no unresolved reference or citation is produced. Printed slips kept literal: no printed word, symbol, hypothesis or number of the counted statement or of its proof was altered. Layout and class: the article's own class is \texttt{siamart220329} with packages including appendix, titletoc, amsfonts, amsmath, fancyhdr, titlesec, indentfirst, booktabs, verbatim, color, bm, xcolor, listings, caption; the wrapper is the lane's pinned \texttt{amsart} editorial wrapper at 10pt on A4, which restates the theorem-like environments the retained text needs and adds the article's own macro definitions that the retained text needs (\texttt{\textbackslash bbx}, \texttt{\textbackslash rd}), with extra wrapper packages (bm, bbm, dsfont, xspace, cancel, mathtools). The delivered numbering is the unit's own. Assets: no retained span contains a figure environment, a graphics-inclusion command, a PDF-page inclusion, a table or a TikZ picture; no image file is copied into this package, the package declares no asset dependency and no image file is redistributed with it. Licence: version arXiv:2407.00533v1 of this article is distributed under the Creative Commons Attribution 4.0 International licence, as recorded in the arXiv licence metadata for this exact version and shown on the abstract page https://arxiv.org/abs/2407.00533v1. A full rights notice is delivered at licenses/arxiv-2407.00533-CC-BY-4.0.txt. Characters: every retained excerpt is pure ASCII, so the delivered engine, XeLaTeX with its default Latin Modern text font, reports no missing character.
\begin{equation} \label{fparticle}
f(t,\bbx) \approx f^N(t,\bbx) := \sum_{p=1}^N w_p \delta (\bbx -\bbx_p(t)),
\end{equation}

\begin{equation} \label{eq:regularE}
E[f](t) = \int_{\mathbb{R}^d} F(t,\bm x, (f*u)(t,\bm x)) + f(t,\bm x)G(t,\bm x,(f*v)(t,\bm x)) \rd \bm x,
\end{equation}

\begin{lemma}
    Let $u,v$ be fixed symmetric integrable functions on $\mathbb{R}^d$ with integrable $\nabla u, \nabla v$. Also, let $F(t,\bm x, y)$, $G(t,\bm x, y)$, $F_y(t,\bm x, y)$, $G_y(t,\bm x, y)$ be integrable functions on $(0,\infty)\times\mathbb{R}^d\times \mathbb{R}$. The energy functional defined in \eqref{eq:regularE} satisfies 
    \begin{equation}
        \frac{1}{w_p}\nabla_{\bm x_p} E[f^N](\bm x_1,\dots,\bm x_N) = \nabla_{\bm x}\frac{\delta E}{\delta f}[f^N](\bm x_p),
    \end{equation} \label{lem:compCond}
 where $f^N$ is the particle approximation given in \eqref{fparticle}.
\end{lemma}

\begin{proof}
For brevity and clarity, we suppress writing dependence on $t$ in $\bm x_q=\bm x_q(t)$. As $(f^N*u)(t,\bm x)=\displaystyle\sum_{q=1}^N w_q u(\bm x-\bm x_q)$ and similarly for $f^N*v$, evaluating $E[f^N]$ yields
    \begin{align*}
        E[f^N] &= \int_{\mathbb{R}^d} F\left(t,\bm x, \sum_{q=1}^N w_q u(\bm x-\bm x_q)\right) \rd\bm x+\sum_{q=1}^N w_q G\left(t,\bm x_q, \sum_{r=1}^N w_r v(\bm x_q-\bm x_r)\right).
    \end{align*} Taking the gradient with respect to $\bm x_p$ gives
        \begin{align}
        \label{eq:LHS-compCond}&\nabla_{\bm x_p} E[f^N](\bm x_1,\dots,\bm x_N)\\
        &=\int_{\mathbb{R}^d} \nabla_{\bm x_p} F\left(t,\bm x, \sum_{q=1}^N w_q u(\bm x-\bm x_q)\right)\rd\bm x \nonumber\\
        &\hskip 2mm+ \nabla_{\bm x_p}\left[w_p G\left(t,\bm x_p,\sum_{r=1}^N w_r v(\bm x_p-\bm x_r)\right)+\sum_{\substack{q=1\\ q\neq p}}^N w_q G\left(t,\bm x_q,\sum_{r=1}^N w_r v(\bm x_q-\bm x_r)\right)\right] \nonumber\\
        &=w_p\int_{\mathbb{R}^d} F_y\left(t,\bm x, \sum_{q=1}^N w_q u(\bm x-\bm x_q)\right) \nabla_{\bbx} u(\bm x_p-\bm x)\rd\bm x\nonumber\\
        &\hskip 2mm+ w_p\left[\nabla_{\bbx}G\left(t,\bm x_p,\sum_{r=1}^N w_r v(\bm x_p-\bm x_r)\right) +G_y\left(t,\bm x_p,\sum_{r=1}^N w_r v(\bm x_p-\bm x_r)\right) \right.\nonumber\\
        &\hskip 6mm \left.\times\sum_{\substack{r=1\\r\neq p}}^N w_r \nabla_{\bbx} v(\bm x_p-\bm x_r) +\sum_{\substack{q=1\\q\neq p}}^N w_qG_y\left(t,\bm x_q,\sum_{r=1}^N w_r v(\bm x_q-\bm x_r)\right) \nabla_{\bbx} v(\bm x_p-\bm x_q)\right].
        \nonumber
    \end{align}
    On the other hand, evaluating the variational derivative of $E[f]$ gives  \begin{align*}
       &\left.\frac{\rd}{\rd\alpha} E[f+\alpha \eta]\right|_{\alpha =0}\\
        &= \int_{\mathbb{R}^d} F_y\left(t, \bm x, (f*u)(t,\bm x)\right)(\eta*u)(\bm x) \rd\bm x\\
        &\hskip 0mm+\int_{\mathbb{R}^d} f(t,\bm x)G_y\left(t, \bm x, (f*v)(t,\bm x)\right)(\eta*v)(\bm x) \rd\bm x + \int_{\mathbb{R}^d} \eta(\bm x)G\left(t,\bm x,(f*v)(t,\bm x)\right) \rd\bm x\\
        &= \int_{\mathbb{R}^d} \eta(\bm y)\dfrac{\delta E}{\delta f}(t,\bm y) \rd\bm y,
    \end{align*}
    where $\displaystyle \frac{\delta E}{\delta f}=F_y(t,\cdot,f*u)*u+\left(fG_y(t,\cdot,f*v)\right)*v+G(t,\cdot,f*v)$.
    Computing the gradient with respect to $\bm x$ of the variational derivative of $E$ yields\begin{align*}
        \nabla_{\bm x}\frac{\delta E}{\delta f}[f](\bm x) &= \int_{\mathbb{R}^d} F_y\left(t, \bm y, (f*u)(t,\bm y)\right)\nabla_{\bbx}u(\bm x-\bm y) \rd\bm y \\
        &\hskip 4mm+ \int_{\mathbb{R}^d} f(t,\bm y)G_y(t,\bm y,(f*v)(t,\bm y))\nabla_{\bbx} v(\bm x-\bm y)\rd\bm y\\
        &\hskip 4mm+ \nabla_{\bbx} G(t,\bm x,(f*v)(\bm x))+G_y(t,\bm x,(f*v)(\bm x))(f*\nabla_{\bbx} v)(\bm x).
    \end{align*}Thus, evaluating $\nabla_{\bm x}\dfrac{\delta E}{\delta f}[f^N]$ at $\bm x=\bm x_p$ gives
    \begin{align}
        \label{eq:RHS-compCond}
        \nabla_{\bm x}\frac{\delta E}{\delta f}[f^N](\bm x_p) &= \int_{\mathbb{R}^d} F_y\left(t, \bm y, \sum_{q=1}^N w_q u(\bm y-\bm x_q)\right)\nabla_{\bbx}u(\bm x_p-\bm y) \rd\bm y\\
        &\hskip 4mm+ \sum_{q=1}^N w_q G_y\left(t,\bm x_q,\sum_{r=1}^N w_r v(\bm x_q-\bm x_r)\right)\nabla_{\bbx} v(\bm x_p-\bm x_q) \nonumber\\
        &\hskip 4mm+\nabla_{\bbx}G\left(t,\bm x_p,\sum_{q=1}^N w_q v(\bm x_p-\bm x_q)\right)\nonumber\\
        &\hskip 4mm+ G_y\left(t,\bm x_p,\sum_{q=1}^N w_q v(\bm x_p-\bm x_q)\right)\left(\sum_{r=1}^N w_r \nabla_{\bbx} v(\bm x_p-\bm x_r)\right).\nonumber
    \end{align}
 Comparing both sides of \eqref{eq:LHS-compCond} and \eqref{eq:RHS-compCond} and using  $\nabla_{\bbx} v(\bm 0) = \bm 0$ (by symmetry of $v$) shows the desired result.
\end{proof}

\end{document}
